% Options for packages loaded elsewhere
% Options for packages loaded elsewhere
\PassOptionsToPackage{unicode}{hyperref}
\PassOptionsToPackage{hyphens}{url}
\PassOptionsToPackage{dvipsnames,svgnames,x11names}{xcolor}
%
\documentclass[
  letterpaper,
  DIV=11,
  numbers=noendperiod]{scrartcl}
\usepackage{xcolor}
\usepackage{amsmath,amssymb}
\setcounter{secnumdepth}{5}
\usepackage{iftex}
\ifPDFTeX
  \usepackage[T1]{fontenc}
  \usepackage[utf8]{inputenc}
  \usepackage{textcomp} % provide euro and other symbols
\else % if luatex or xetex
  \usepackage{unicode-math} % this also loads fontspec
  \defaultfontfeatures{Scale=MatchLowercase}
  \defaultfontfeatures[\rmfamily]{Ligatures=TeX,Scale=1}
\fi
\usepackage{lmodern}
\ifPDFTeX\else
  % xetex/luatex font selection
  \setmainfont[]{Times New Roman}
\fi
% Use upquote if available, for straight quotes in verbatim environments
\IfFileExists{upquote.sty}{\usepackage{upquote}}{}
\IfFileExists{microtype.sty}{% use microtype if available
  \usepackage[]{microtype}
  \UseMicrotypeSet[protrusion]{basicmath} % disable protrusion for tt fonts
}{}
\makeatletter
\@ifundefined{KOMAClassName}{% if non-KOMA class
  \IfFileExists{parskip.sty}{%
    \usepackage{parskip}
  }{% else
    \setlength{\parindent}{0pt}
    \setlength{\parskip}{6pt plus 2pt minus 1pt}}
}{% if KOMA class
  \KOMAoptions{parskip=half}}
\makeatother
% Make \paragraph and \subparagraph free-standing
\makeatletter
\ifx\paragraph\undefined\else
  \let\oldparagraph\paragraph
  \renewcommand{\paragraph}{
    \@ifstar
      \xxxParagraphStar
      \xxxParagraphNoStar
  }
  \newcommand{\xxxParagraphStar}[1]{\oldparagraph*{#1}\mbox{}}
  \newcommand{\xxxParagraphNoStar}[1]{\oldparagraph{#1}\mbox{}}
\fi
\ifx\subparagraph\undefined\else
  \let\oldsubparagraph\subparagraph
  \renewcommand{\subparagraph}{
    \@ifstar
      \xxxSubParagraphStar
      \xxxSubParagraphNoStar
  }
  \newcommand{\xxxSubParagraphStar}[1]{\oldsubparagraph*{#1}\mbox{}}
  \newcommand{\xxxSubParagraphNoStar}[1]{\oldsubparagraph{#1}\mbox{}}
\fi
\makeatother


\usepackage{longtable,booktabs,array}
\newcounter{none} % for unnumbered tables
\usepackage{calc} % for calculating minipage widths
% Correct order of tables after \paragraph or \subparagraph
\usepackage{etoolbox}
\makeatletter
\patchcmd\longtable{\par}{\if@noskipsec\mbox{}\fi\par}{}{}
\makeatother
% Allow footnotes in longtable head/foot
\IfFileExists{footnotehyper.sty}{\usepackage{footnotehyper}}{\usepackage{footnote}}
\makesavenoteenv{longtable}
\usepackage{graphicx}
\makeatletter
\newsavebox\pandoc@box
\newcommand*\pandocbounded[1]{% scales image to fit in text height/width
  \sbox\pandoc@box{#1}%
  \Gscale@div\@tempa{\textheight}{\dimexpr\ht\pandoc@box+\dp\pandoc@box\relax}%
  \Gscale@div\@tempb{\linewidth}{\wd\pandoc@box}%
  \ifdim\@tempb\p@<\@tempa\p@\let\@tempa\@tempb\fi% select the smaller of both
  \ifdim\@tempa\p@<\p@\scalebox{\@tempa}{\usebox\pandoc@box}%
  \else\usebox{\pandoc@box}%
  \fi%
}
% Set default figure placement to htbp
\def\fps@figure{htbp}
\makeatother





\setlength{\emergencystretch}{3em} % prevent overfull lines

\providecommand{\tightlist}{%
  \setlength{\itemsep}{0pt}\setlength{\parskip}{0pt}}



 


\KOMAoption{captions}{tableheading}
\makeatletter
\@ifpackageloaded{caption}{}{\usepackage{caption}}
\AtBeginDocument{%
\ifdefined\contentsname
  \renewcommand*\contentsname{Table of contents}
\else
  \newcommand\contentsname{Table of contents}
\fi
\ifdefined\listfigurename
  \renewcommand*\listfigurename{List of Figures}
\else
  \newcommand\listfigurename{List of Figures}
\fi
\ifdefined\listtablename
  \renewcommand*\listtablename{List of Tables}
\else
  \newcommand\listtablename{List of Tables}
\fi
\ifdefined\figurename
  \renewcommand*\figurename{Figure}
\else
  \newcommand\figurename{Figure}
\fi
\ifdefined\tablename
  \renewcommand*\tablename{Table}
\else
  \newcommand\tablename{Table}
\fi
}
\@ifpackageloaded{float}{}{\usepackage{float}}
\floatstyle{ruled}
\@ifundefined{c@chapter}{\newfloat{codelisting}{h}{lop}}{\newfloat{codelisting}{h}{lop}[chapter]}
\floatname{codelisting}{Listing}
\newcommand*\listoflistings{\listof{codelisting}{List of Listings}}
\makeatother
\makeatletter
\makeatother
\makeatletter
\@ifpackageloaded{caption}{}{\usepackage{caption}}
\@ifpackageloaded{subcaption}{}{\usepackage{subcaption}}
\makeatother
\usepackage{bookmark}
\IfFileExists{xurl.sty}{\usepackage{xurl}}{} % add URL line breaks if available
\urlstyle{same}
\hypersetup{
  pdftitle={CV exploration: is the path floor RMSE-optimal{,} and can a selector beat it?},
  colorlinks=true,
  linkcolor={blue},
  filecolor={Maroon},
  citecolor={Blue},
  urlcolor={Blue},
  pdfcreator={LaTeX via pandoc}}


\title{CV exploration: is the path floor RMSE-optimal, and can a
selector beat it?}
\author{}
\date{}
\begin{document}
\maketitle


\section{Investigation}\label{investigation}

2 aims:

\textbf{A:} Is the optimal (minimizes RMSE) \(\lambda\) interior,
i.e.~not the path floor \(\lambda_{\min}\)?\\
\textbf{B:} If so, does any selector beat the path floor?

\textbf{Target criterion}

All selectors use \texttt{X} and \texttt{W} only and return
\texttt{lambda.min}, the path value minimising their criterion on
held-out units (treated arm):

\begin{itemize}
\tightlist
\item
  \texttt{bloss}: balance loss \(\frac{1}{n}\sum_i [W_i e^{-\eta_i} +
  (1 - W_i)\eta_i]\), the balnet objective (Tan, 2020).
\item
  \texttt{smd}: mean \(|\mathrm{SMD}_j|\), IPW-weighted treated mean
  against the full-sample mean, in full-sample SDs.
\item
  \texttt{inf}: max \(|\mathrm{SMD}_j|\).
\end{itemize}

\textbf{CV vs bootstrap}

\begin{itemize}
\tightlist
\item
  CV (\texttt{cv.balnet}, 10 folds): refit on the training folds,
  criterion on the held-out fold, averaged over folds.
\item
  Bootstrap (\texttt{cv.boot.balnet}, 500 half-samples): no refitting;
  the full-sample weights are scored on random half-samples against the
  full-sample means and SDs (Wang and Zubizarreta, 2020).
\end{itemize}

\textbf{notes:}

\begin{itemize}
\tightlist
\item
  \emph{Path floor} \(\lambda_{\min}\) \emph{is balnet default,}
  \(\lambda_{\max}/100\)\emph{, stated if changed.}
\item
  \emph{Sampling noise: 500 reps, relative MCSE 3\% (Morris et al,
  2019).}
\item
  \emph{Default penalty} \(\alpha = 1\)\emph{, lasso regularisation
  (stated if changed).}
\end{itemize}

\section{Outcome Noise, Density, and
Overlap}\label{outcome-noise-density-and-overlap}

\subsection{Sparse outcome}\label{sparse-outcome}

\begin{itemize}
\tightlist
\item
  5 active confounders (\texttt{s\_y\ =\ 5}).
\item
  \texttt{sigma\_y} increases outcome noise (reduces \(R^2\), lowers
  SNR).
\item
  \texttt{c\_prop} \(\in \{0.7, 1.5, 2.5\}\) for \texttt{overlap} good
  \(>\) moderate \(>\) bad: degrades overlap.
\end{itemize}

\includegraphics[width=1\linewidth,height=\textheight,keepaspectratio]{fig_density_sy5.png}

\textbf{A:} with \texttt{sigma\_y\ =\ 5}, interior solutions exist at
all 3 overlap levels.

\textbf{B:} With large outcome noise, selectors outperform the path
floor \(\lambda_{\min}\) at good and moderate overlap; at bad overlap
the best selector only ties it.

\begin{itemize}
\item
  good overlap: \texttt{boot.smd} \(\approx\) \texttt{boot.inf}
  \textgreater{} \texttt{cv.bloss} \textgreater{} path floor
  \(\lambda_{\min}\)
\item
  moderate overlap: \texttt{boot.inf} \textgreater{} \texttt{cv.bloss}
  \textgreater{} path floor \(\lambda_{\min}\)
\item
  bad overlap: \texttt{cv.bloss} \(\approx\) path floor
  \(\lambda_{\min}\)
\end{itemize}

\textbf{Increasing outcome noise causes both interior optimums and
selector wins.}

\textbf{Criterion selection:} Degrading overlap increases the relative
performance of the path floor \(\lambda_{\min}\). It also changes the
ranking between selection criteria (\texttt{boot.smd} collapses beyond
good overlap, \texttt{boot.inf} holds, \texttt{cv.bloss} improves
relative to both SMD criteria).

\textbf{Bootstrapping vs cross-validation:} for both criteria where the
two are available (\texttt{smd}, \texttt{inf}), bootstrapping
outperforms CV at every overlap level. \texttt{cv.bloss} is the only
competitive criterion under bad overlap, and has no bootstrapped
equivalent.

\emph{Note: sampling error affects ranking (Appendix 2). Same DGP, same
reps: at bad overlap, \texttt{sigma\_y\ =\ 5}, \texttt{cv.bloss} now
edges the path floor and \texttt{boot.inf} ties it.}

\subsection{Dense outcome: fixed
confounding}\label{dense-outcome-fixed-confounding}

DGP isolates the effect of outcome density by holding confounder signal
fixed:

\begin{itemize}
\tightlist
\item
  50 active outcome covariates:

  \begin{itemize}
  \tightlist
  \item
    5 confounders at \texttt{1/sqrt(5)}
  \item
    45 outcome-only predictors at \texttt{1/sqrt(45)}
  \end{itemize}
\end{itemize}

Appendix 1 shows a diluted case where
\texttt{beta\_y{[}1:s\_y{]}\ \textless{}-\ 1/sqrt(s\_y)} makes
increasing density also decrease confounder signal (in effect,
increasing \texttt{sigma\_y}).

\includegraphics[width=1\linewidth,height=\textheight,keepaspectratio]{fig_density_fixed_sigma5.png}

Under both high and low noise settings \emph{(Appendix 3)} outcome
density does not appear to drive selection performance, or interior
minimums. Results for density show only minor changes to relative
performance, with any ranking differences likely sampling noise.

\subsection{Investigating noise X
overlap.}\label{investigating-noise-x-overlap.}

\includegraphics[width=23.2in,height=\textheight,keepaspectratio]{fig_snr_lasso.png}

\emph{cell (5,3 - sigma 10 bad overlap at 1000 reps to confirm pattern
change.}

\textbf{A:} interior minimums appear at lower noise under strong
overlap; as overlap degrades, more noise is needed.

\begin{itemize}
\tightlist
\item
  good overlap: interior at \(\sigma_y = 3\) (possibly 2, a draw by eye)
\item
  moderate overlap: interior from \(\sigma_y = 3\) (faint at 3, clear at
  5)
\item
  bad overlap: interior at \(\sigma_y = 5\)
\end{itemize}

\textbf{B:} Selectors beat the path floor \(\lambda_{\min}\) once the
interior is clear. Best selector changes with noise and overlap.

\begin{itemize}
\tightlist
\item
  good overlap:

  \begin{itemize}
  \tightlist
  \item
    \(\sigma_y = 2\): path floor \(\lambda_{\min}\) \(\approx\)
    \texttt{boot.inf} \(\approx\) \texttt{boot.smd} \emph{(draw)}
  \item
    \(\sigma_y = 3\): \texttt{boot.inf} \(\approx\) \texttt{boot.smd}
    \textgreater{} path floor \(\lambda_{\min}\)
  \item
    \(\sigma_y = 5\): \texttt{boot.inf} \(\approx\) \texttt{boot.smd}
    \textgreater{} \texttt{cv.bloss} \(\approx\) path floor
    \(\lambda_{\min}\)
  \item
    \(\sigma_y = 10\): \texttt{cv.bloss} \textgreater{}
    \texttt{boot.inf} \(\approx\) \texttt{boot.smd} \textgreater{}
    \texttt{cv.inf} \(\approx\) path floor \(\lambda_{\min}\)
  \end{itemize}
\item
  moderate overlap:

  \begin{itemize}
  \tightlist
  \item
    \(\sigma_y = 5\): \texttt{boot.inf} \textgreater{} \texttt{cv.bloss}
    \(\approx\) path floor \(\lambda_{\min}\)
  \item
    \(\sigma_y = 10\): \texttt{cv.bloss} \textgreater{}
    \texttt{boot.inf} \textgreater{} \texttt{boot.smd} \(\approx\)
    \texttt{cv.inf} \textgreater{} path floor \(\lambda_{\min}\)
  \end{itemize}
\item
  bad overlap:

  \begin{itemize}
  \tightlist
  \item
    \(\sigma_y = 5\): \texttt{cv.bloss} \(\approx\) path floor
    \(\lambda_{\min}\) (draw)
  \item
    \(\sigma_y = 10\): \texttt{boot.inf} \textgreater{}
    \texttt{cv.bloss} \(\approx\) \texttt{cv.inf} \textgreater{}
    \texttt{boot.smd} \textgreater{} path floor \(\lambda_{\min}\)
  \end{itemize}
\end{itemize}

\textbf{Interior minimums:} Interior minimums depend on both noise and
overlap. \(\lambda_{\min}\) is more robust to degredation in overlap,
while the CV/BS selectors that are competitive, are generally more
robust to increased outcome noise.

\textbf{Criterion selection:} \texttt{cv.bloss} and \texttt{boot.inf}
are the most effective criteria. \texttt{cv.bloss} is more robust to
overlap degradation, outperforming \texttt{boot.inf} under ``bad''
overlap for \(\sigma_y \le 5\) and underperforming for ``good'' and
``moderate'' overlap.

This pattern inverses between \(\sigma_y = 7\) and \(\sigma_y = 10\),
where the order flips with \texttt{cv.bloss} \textgreater{}
\texttt{boot.inf} for good and moderate overlap, and \texttt{cv.bloss}
\textless{} \texttt{boot.inf} for bad overlap. \textbf{Why?}

\begin{itemize}
\tightlist
\item
  \emph{Possibly:} nas either selector observed outcome, each pick is
  fixed while noise pushes the optimal \(\lambda\) up. The deeper pick
  wins at low \(\sigma_y\), the earlier at high \(\sigma_y\), and
  overlap sets which selector stops first: \texttt{boot.inf} stops
  deeper under good and moderate overlap, earlier under bad.
\end{itemize}

\emph{Holds at 1000 reps, explored below.}

\subsection{Independent: noise and
overlap}\label{independent-noise-and-overlap}

\begin{figure}[p]

\centering{

\includegraphics[width=15in,height=\textheight,keepaspectratio]{fig_axes_noise_overlap.png}

}

\caption{\label{fig-axes}}

\end{figure}%

\textbf{Figure:} LHS: increasing noise, (sigma\_y=1 to 18). RHS:
degrading overlap(0.7 to 16), sigma\_y=1.

\textbf{Overlap and Noise, Seperatly:}

\begin{itemize}
\item
  \textbf{overlap axis} (\(\sigma_y = 1\)): there is no interior
  solution, the path is still falling, indicating \(\lambda_{\min}\) is
  best at every level and its apparent losses in the figure
  (\texttt{c\_prop} 6 to 16) are likely artifacts . The order of
  \texttt{cv.bloss} and \texttt{boot.inf} flips between \texttt{c\_prop}
  2 and 2.5.

  \begin{itemize}
  \tightlist
  \item
    \texttt{c\_prop} = 0.7 to 1.5: \texttt{boot.inf} \textgreater{}
    \texttt{cv.bloss}
  \item
    \texttt{c\_prop} = 2: \texttt{cv.bloss} \(\approx\)
    \texttt{boot.inf}
  \item
    \texttt{c\_prop} = 2.5 to 16: \texttt{cv.bloss} \textgreater{}
    \texttt{boot.inf}
  \end{itemize}
\item
  \textbf{noise axis} (good overlap): \(\lambda_{\min}\) is best only at
  \(\sigma_y\) = 1; \texttt{boot.inf} ties it at 2 and beats it from 3,
  \texttt{cv.bloss} ties it at 5 and beats it from 7; the order of
  \texttt{cv.bloss} and \texttt{boot.inf} flips between \(\sigma_y\) 7
  and 10.

  \begin{itemize}
  \item
    \(\sigma_y\) = 1: \(\lambda_{\min}\) \textgreater{}
    \texttt{boot.inf} \textgreater{} \texttt{cv.bloss}
  \item
    \(\sigma_y\) = 2: \texttt{boot.inf} \(\approx\) \(\lambda_{\min}\)
    \textgreater{} \texttt{cv.bloss}
  \item
    \(\sigma_y\) = 3 to 4: \texttt{boot.inf} \textgreater{}
    \(\lambda_{\min}\) \textgreater{} \texttt{cv.bloss}
  \item
    \(\sigma_y\) = 5: \texttt{boot.inf} \textgreater{} \texttt{cv.bloss}
    \(\approx\) \(\lambda_{\min}\)
  \item
    \(\sigma_y\) = 7: \texttt{cv.bloss} \(\approx\) \texttt{boot.inf}
    \textgreater{} \(\lambda_{\min}\)
  \item
    \(\sigma_y\) = 10 to 18: \texttt{cv.bloss} \textgreater{}
    \texttt{boot.inf} \textgreater{} \(\lambda_{\min}\)
  \end{itemize}
\end{itemize}

\textbf{New evidence for bootstrapping vs cross-validation:} earlier
finding (bootstrap \textgreater{} CV for \texttt{smd}, \texttt{inf} at
every overlap level) holds on the overlap axis, breaks on the noise
axis.

\begin{itemize}
\tightlist
\item
  overlap axis: bootstrap \textgreater{} CV at every \texttt{c\_prop}
\item
  noise axis: bootstrapping's relative advantage fades when noise is
  high enough.

  \begin{itemize}
  \tightlist
  \item
    \texttt{cv.inf} \textgreater{} \texttt{boot.inf} from \(\sigma_y\) =
    12
  \item
    \texttt{cv.smd} \textgreater{} \texttt{boot.smd} by 18
  \end{itemize}
\end{itemize}

\textbf{The pattern and the flip:} our findings support the initial
pattern as:

\begin{itemize}
\tightlist
\item
  At good overlap, \texttt{boot.inf} \textgreater{} \texttt{cv.bloss},
  by \texttt{c\_prop\ =\ 2} they perform the same, by bad overlap,
  \texttt{cv.bloss} \textgreater{} \texttt{boot.inf}.
\item
  \texttt{boot.inf} \textgreater{} \texttt{cv.bloss} for noise up to 5,
  at noise 7 \texttt{cv.bloss} \(\approx\) \texttt{boot.inf}, by 10
  \texttt{cv.bloss} \textgreater{} \texttt{boot.inf}.
\end{itemize}

These are consistent with the initial pattern. The only change in
relative performance we see on the noise axis is that \texttt{cv.bloss}
outperforms \texttt{boot.inf} somewhere over \(\sigma_y\) = 7. This is
consistent with its relative performance overtaking \texttt{boot.inf} at
\(\sigma_y\) = 10 for good and moderate overlap. \textbf{the flip under
bad overlap is not supported by this evidence}., where \texttt{boot.inf}
\textgreater{} \texttt{cv.bloss} \(\sigma_y\) = 10, is not supported
however, \textbf{Why?}

\begin{figure}

\begin{minipage}[t]{0.50\linewidth}

{\def\LTcaptype{none} % do not increment counter
\begin{longtable}[]{@{}rrrr@{}}
\toprule\noalign{}
sigma\_y & lambda\_opt & cv.bloss & boot.inf \\
\midrule\noalign{}
\endhead
\bottomrule\noalign{}
\endlastfoot
1 & 0.005 & 0.069 & 0.017 \\
2 & 0.011 & 0.068 & 0.017 \\
3 & 0.017 & 0.067 & 0.017 \\
4 & 0.018 & 0.067 & 0.017 \\
5 & 0.032 & 0.067 & 0.017 \\
7 & 0.048 & 0.068 & 0.017 \\
10 & 0.058 & 0.068 & 0.017 \\
12 & 0.067 & 0.068 & 0.017 \\
14 & 0.070 & 0.068 & 0.017 \\
16 & 0.085 & 0.067 & 0.017 \\
18 & 0.092 & 0.068 & 0.017 \\
20 & 0.102 & 0.069 & 0.017 \\
\end{longtable}
}

\end{minipage}%
%
\begin{minipage}[t]{0.50\linewidth}

{\def\LTcaptype{none} % do not increment counter
\begin{longtable}[]{@{}rrrr@{}}
\toprule\noalign{}
c\_prop & lambda\_opt & cv.bloss & boot.inf \\
\midrule\noalign{}
\endhead
\bottomrule\noalign{}
\endlastfoot
0.7 & 0.005 & 0.069 & 0.017 \\
1.0 & 0.014 & 0.079 & 0.027 \\
1.5 & 0.026 & 0.105 & 0.058 \\
2.0 & 0.061 & 0.132 & 0.143 \\
2.5 & 0.092 & 0.153 & 0.191 \\
3.0 & 0.116 & 0.169 & 0.211 \\
4.0 & 0.168 & 0.198 & 0.238 \\
5.0 & 0.199 & 0.218 & 0.256 \\
6.0 & 0.185 & 0.229 & 0.266 \\
8.0 & 0.242 & 0.247 & 0.278 \\
16.0 & 0.292 & 0.265 & 0.293 \\
\end{longtable}
}

\end{minipage}%

\end{figure}%

As selection models are outcome blind, increasing noise doesnt change
their selection lambda, it only raises optimal level of imbalance, CV
bloss selects higher, so when noise is high enough - it wins.

Interestingly, at c\_prop=16, boot inf is closer to lambda optimum but
still has a higher RMSE, **possibly artifact, maybe investigate.**

\begin{longtable}[]{@{}
  >{\raggedleft\arraybackslash}p{(\linewidth - 12\tabcolsep) * \real{0.1231}}
  >{\raggedleft\arraybackslash}p{(\linewidth - 12\tabcolsep) * \real{0.0923}}
  >{\raggedleft\arraybackslash}p{(\linewidth - 12\tabcolsep) * \real{0.1846}}
  >{\raggedleft\arraybackslash}p{(\linewidth - 12\tabcolsep) * \real{0.1385}}
  >{\raggedleft\arraybackslash}p{(\linewidth - 12\tabcolsep) * \real{0.1846}}
  >{\raggedleft\arraybackslash}p{(\linewidth - 12\tabcolsep) * \real{0.0923}}
  >{\raggedleft\arraybackslash}p{(\linewidth - 12\tabcolsep) * \real{0.1846}}@{}}
\caption{RMSE-optimal lambda}\tabularnewline
\toprule\noalign{}
\begin{minipage}[b]{\linewidth}\raggedleft
sigma\_y
\end{minipage} & \begin{minipage}[b]{\linewidth}\raggedleft
good
\end{minipage} & \begin{minipage}[b]{\linewidth}\raggedleft
cum. change
\end{minipage} & \begin{minipage}[b]{\linewidth}\raggedleft
moderate
\end{minipage} & \begin{minipage}[b]{\linewidth}\raggedleft
cum. change
\end{minipage} & \begin{minipage}[b]{\linewidth}\raggedleft
bad
\end{minipage} & \begin{minipage}[b]{\linewidth}\raggedleft
cum. change
\end{minipage} \\
\midrule\noalign{}
\endfirsthead
\toprule\noalign{}
\begin{minipage}[b]{\linewidth}\raggedleft
sigma\_y
\end{minipage} & \begin{minipage}[b]{\linewidth}\raggedleft
good
\end{minipage} & \begin{minipage}[b]{\linewidth}\raggedleft
cum. change
\end{minipage} & \begin{minipage}[b]{\linewidth}\raggedleft
moderate
\end{minipage} & \begin{minipage}[b]{\linewidth}\raggedleft
cum. change
\end{minipage} & \begin{minipage}[b]{\linewidth}\raggedleft
bad
\end{minipage} & \begin{minipage}[b]{\linewidth}\raggedleft
cum. change
\end{minipage} \\
\midrule\noalign{}
\endhead
\bottomrule\noalign{}
\endlastfoot
1 & 0.005 & 1.00 & 0.026 & 1.00 & 0.092 & 1.00 \\
2 & 0.011 & 2.12 & 0.039 & 1.49 & 0.125 & 1.36 \\
3 & 0.017 & 3.42 & 0.050 & 1.88 & 0.090 & 0.98 \\
4 & 0.018 & 3.56 & 0.057 & 2.16 & 0.119 & 1.30 \\
5 & 0.032 & 6.39 & 0.064 & 2.40 & 0.122 & 1.33 \\
7 & 0.048 & 9.44 & 0.081 & 3.08 & 0.173 & 1.89 \\
10 & 0.058 & 11.39 & 0.134 & 5.08 & 0.218 & 2.37 \\
12 & 0.067 & 13.31 & NA & NA & NA & NA \\
\end{longtable}

\begin{longtable}[]{@{}
  >{\raggedleft\arraybackslash}p{(\linewidth - 12\tabcolsep) * \real{0.1231}}
  >{\raggedleft\arraybackslash}p{(\linewidth - 12\tabcolsep) * \real{0.0923}}
  >{\raggedleft\arraybackslash}p{(\linewidth - 12\tabcolsep) * \real{0.1846}}
  >{\raggedleft\arraybackslash}p{(\linewidth - 12\tabcolsep) * \real{0.1385}}
  >{\raggedleft\arraybackslash}p{(\linewidth - 12\tabcolsep) * \real{0.1846}}
  >{\raggedleft\arraybackslash}p{(\linewidth - 12\tabcolsep) * \real{0.0923}}
  >{\raggedleft\arraybackslash}p{(\linewidth - 12\tabcolsep) * \real{0.1846}}@{}}
\caption{cv.bloss median pick}\tabularnewline
\toprule\noalign{}
\begin{minipage}[b]{\linewidth}\raggedleft
sigma\_y
\end{minipage} & \begin{minipage}[b]{\linewidth}\raggedleft
good
\end{minipage} & \begin{minipage}[b]{\linewidth}\raggedleft
cum. change
\end{minipage} & \begin{minipage}[b]{\linewidth}\raggedleft
moderate
\end{minipage} & \begin{minipage}[b]{\linewidth}\raggedleft
cum. change
\end{minipage} & \begin{minipage}[b]{\linewidth}\raggedleft
bad
\end{minipage} & \begin{minipage}[b]{\linewidth}\raggedleft
cum. change
\end{minipage} \\
\midrule\noalign{}
\endfirsthead
\toprule\noalign{}
\begin{minipage}[b]{\linewidth}\raggedleft
sigma\_y
\end{minipage} & \begin{minipage}[b]{\linewidth}\raggedleft
good
\end{minipage} & \begin{minipage}[b]{\linewidth}\raggedleft
cum. change
\end{minipage} & \begin{minipage}[b]{\linewidth}\raggedleft
moderate
\end{minipage} & \begin{minipage}[b]{\linewidth}\raggedleft
cum. change
\end{minipage} & \begin{minipage}[b]{\linewidth}\raggedleft
bad
\end{minipage} & \begin{minipage}[b]{\linewidth}\raggedleft
cum. change
\end{minipage} \\
\midrule\noalign{}
\endhead
\bottomrule\noalign{}
\endlastfoot
1 & 0.069 & 1.00 & 0.105 & 1.00 & 0.153 & 1.00 \\
2 & 0.068 & 0.99 & 0.106 & 1.01 & 0.153 & 1.00 \\
3 & 0.067 & 0.98 & 0.106 & 1.01 & 0.154 & 1.00 \\
4 & 0.067 & 0.97 & 0.107 & 1.02 & 0.153 & 1.00 \\
5 & 0.067 & 0.98 & 0.106 & 1.01 & 0.152 & 0.99 \\
7 & 0.068 & 0.99 & 0.106 & 1.01 & 0.154 & 1.00 \\
10 & 0.068 & 0.98 & 0.104 & 0.99 & 0.152 & 0.99 \\
12 & 0.068 & 0.98 & NA & NA & NA & NA \\
\end{longtable}

\begin{longtable}[]{@{}
  >{\raggedleft\arraybackslash}p{(\linewidth - 12\tabcolsep) * \real{0.1231}}
  >{\raggedleft\arraybackslash}p{(\linewidth - 12\tabcolsep) * \real{0.0923}}
  >{\raggedleft\arraybackslash}p{(\linewidth - 12\tabcolsep) * \real{0.1846}}
  >{\raggedleft\arraybackslash}p{(\linewidth - 12\tabcolsep) * \real{0.1385}}
  >{\raggedleft\arraybackslash}p{(\linewidth - 12\tabcolsep) * \real{0.1846}}
  >{\raggedleft\arraybackslash}p{(\linewidth - 12\tabcolsep) * \real{0.0923}}
  >{\raggedleft\arraybackslash}p{(\linewidth - 12\tabcolsep) * \real{0.1846}}@{}}
\caption{boot.inf median pick}\tabularnewline
\toprule\noalign{}
\begin{minipage}[b]{\linewidth}\raggedleft
sigma\_y
\end{minipage} & \begin{minipage}[b]{\linewidth}\raggedleft
good
\end{minipage} & \begin{minipage}[b]{\linewidth}\raggedleft
cum. change
\end{minipage} & \begin{minipage}[b]{\linewidth}\raggedleft
moderate
\end{minipage} & \begin{minipage}[b]{\linewidth}\raggedleft
cum. change
\end{minipage} & \begin{minipage}[b]{\linewidth}\raggedleft
bad
\end{minipage} & \begin{minipage}[b]{\linewidth}\raggedleft
cum. change
\end{minipage} \\
\midrule\noalign{}
\endfirsthead
\toprule\noalign{}
\begin{minipage}[b]{\linewidth}\raggedleft
sigma\_y
\end{minipage} & \begin{minipage}[b]{\linewidth}\raggedleft
good
\end{minipage} & \begin{minipage}[b]{\linewidth}\raggedleft
cum. change
\end{minipage} & \begin{minipage}[b]{\linewidth}\raggedleft
moderate
\end{minipage} & \begin{minipage}[b]{\linewidth}\raggedleft
cum. change
\end{minipage} & \begin{minipage}[b]{\linewidth}\raggedleft
bad
\end{minipage} & \begin{minipage}[b]{\linewidth}\raggedleft
cum. change
\end{minipage} \\
\midrule\noalign{}
\endhead
\bottomrule\noalign{}
\endlastfoot
1 & 0.017 & 1.00 & 0.058 & 1.00 & 0.191 & 1.00 \\
2 & 0.017 & 1.00 & 0.057 & 0.98 & 0.187 & 0.98 \\
3 & 0.017 & 1.01 & 0.057 & 0.98 & 0.190 & 1.00 \\
4 & 0.017 & 0.98 & 0.056 & 0.97 & 0.187 & 0.98 \\
5 & 0.017 & 0.97 & 0.057 & 0.99 & 0.186 & 0.98 \\
7 & 0.017 & 0.99 & 0.056 & 0.96 & 0.188 & 0.99 \\
10 & 0.017 & 0.98 & 0.056 & 0.97 & 0.189 & 0.99 \\
12 & 0.017 & 0.99 & NA & NA & NA & NA \\
\end{longtable}

\begin{longtable}[]{@{}
  >{\raggedleft\arraybackslash}p{(\linewidth - 12\tabcolsep) * \real{0.1231}}
  >{\raggedleft\arraybackslash}p{(\linewidth - 12\tabcolsep) * \real{0.0923}}
  >{\raggedleft\arraybackslash}p{(\linewidth - 12\tabcolsep) * \real{0.1846}}
  >{\raggedleft\arraybackslash}p{(\linewidth - 12\tabcolsep) * \real{0.1385}}
  >{\raggedleft\arraybackslash}p{(\linewidth - 12\tabcolsep) * \real{0.1846}}
  >{\raggedleft\arraybackslash}p{(\linewidth - 12\tabcolsep) * \real{0.0923}}
  >{\raggedleft\arraybackslash}p{(\linewidth - 12\tabcolsep) * \real{0.1846}}@{}}
\caption{cv.inf median pick}\tabularnewline
\toprule\noalign{}
\begin{minipage}[b]{\linewidth}\raggedleft
sigma\_y
\end{minipage} & \begin{minipage}[b]{\linewidth}\raggedleft
good
\end{minipage} & \begin{minipage}[b]{\linewidth}\raggedleft
cum. change
\end{minipage} & \begin{minipage}[b]{\linewidth}\raggedleft
moderate
\end{minipage} & \begin{minipage}[b]{\linewidth}\raggedleft
cum. change
\end{minipage} & \begin{minipage}[b]{\linewidth}\raggedleft
bad
\end{minipage} & \begin{minipage}[b]{\linewidth}\raggedleft
cum. change
\end{minipage} \\
\midrule\noalign{}
\endfirsthead
\toprule\noalign{}
\begin{minipage}[b]{\linewidth}\raggedleft
sigma\_y
\end{minipage} & \begin{minipage}[b]{\linewidth}\raggedleft
good
\end{minipage} & \begin{minipage}[b]{\linewidth}\raggedleft
cum. change
\end{minipage} & \begin{minipage}[b]{\linewidth}\raggedleft
moderate
\end{minipage} & \begin{minipage}[b]{\linewidth}\raggedleft
cum. change
\end{minipage} & \begin{minipage}[b]{\linewidth}\raggedleft
bad
\end{minipage} & \begin{minipage}[b]{\linewidth}\raggedleft
cum. change
\end{minipage} \\
\midrule\noalign{}
\endhead
\bottomrule\noalign{}
\endlastfoot
1 & 0.149 & 1.00 & 0.241 & 1.00 & 0.302 & 1.00 \\
2 & 0.144 & 0.97 & 0.243 & 1.01 & 0.301 & 1.00 \\
3 & 0.150 & 1.01 & 0.247 & 1.02 & 0.302 & 1.00 \\
4 & 0.148 & 1.00 & 0.242 & 1.00 & 0.302 & 1.00 \\
5 & 0.145 & 0.97 & 0.243 & 1.01 & 0.301 & 1.00 \\
7 & 0.147 & 0.99 & 0.240 & 1.00 & 0.301 & 1.00 \\
10 & 0.145 & 0.97 & 0.243 & 1.01 & 0.300 & 0.99 \\
12 & 0.147 & 0.99 & NA & NA & NA & NA \\
\end{longtable}

\begin{figure}

\centering{

\pandocbounded{\includegraphics[keepaspectratio,alt={Median selected lambda for cv.bloss and boot.inf against the RMSE-optimal lambda, along the noise axis (good overlap) and the overlap axis (\$\textbackslash sigma\_y = 1\$). Log scales.}]{fig_picks_vs_opt.png}}

}

\caption{\label{fig-picks-vs-opt}}

\end{figure}%

\begin{figure}

\centering{

\pandocbounded{\includegraphics[keepaspectratio]{fig_picks_grid.png}}

}

\caption{\label{fig-picks-grid}Median selected lambda for
\texttt{cv.bloss} and \texttt{boot.inf} against the RMSE-optimal lambda,
log scale.}

\end{figure}%

\textbf{Some of the ``how'', maybe some of the why:}

\begin{enumerate}
\def\labelenumi{\arabic{enumi}.}
\tightlist
\item
  \textbf{Noise dynamics.} The cumulative interaction of degraded
  overlap with increased noise, causes the optimal lambda to increase
  outweighing the robustness of targeting balancing loss, (cv or not).
  Both boot.inf and cv.inf will outperform cv.bloss when noise is high
  enough.
\end{enumerate}

\begin{enumerate}
\def\labelenumi{\arabic{enumi}.}
\setcounter{enumi}{1}
\tightlist
\item
  \textbf{Flip dynamic related to criterion not bootstrapping.} When
  outcome isnt visible, the \texttt{inf} (max \(|\mathrm{SMD}_j|\))
  criterion variants, which minimise the worst balanced co-variate,
  always choose a higher lambda, making them more robust to highly
  degraded overlap, since as overlap weakens the optimal RMSE rises.
\end{enumerate}

The bootstrap component is not relevent to the flip, with
cross-validation likely to perform worse when overlap degrades past a
certain point, given.

\begin{itemize}
\tightlist
\item
  \texttt{cv.inf} \textgreater{} \texttt{boot.inf} from \(\sigma_y\) =
  12
\item
  \texttt{cv.smd} \textgreater{} \texttt{boot.smd} by 18
\end{itemize}

\begin{figure}

\begin{minipage}[t]{0.50\linewidth}

{\def\LTcaptype{none} % do not increment counter
\begin{longtable}[]{@{}
  >{\raggedleft\arraybackslash}p{(\linewidth - 12\tabcolsep) * \real{0.0988}}
  >{\raggedleft\arraybackslash}p{(\linewidth - 12\tabcolsep) * \real{0.1358}}
  >{\raggedleft\arraybackslash}p{(\linewidth - 12\tabcolsep) * \real{0.1481}}
  >{\raggedleft\arraybackslash}p{(\linewidth - 12\tabcolsep) * \real{0.1481}}
  >{\raggedleft\arraybackslash}p{(\linewidth - 12\tabcolsep) * \real{0.1728}}
  >{\raggedleft\arraybackslash}p{(\linewidth - 12\tabcolsep) * \real{0.1358}}
  >{\raggedleft\arraybackslash}p{(\linewidth - 12\tabcolsep) * \real{0.1605}}@{}}
\toprule\noalign{}
\begin{minipage}[b]{\linewidth}\raggedleft
sigma\_y
\end{minipage} & \begin{minipage}[b]{\linewidth}\raggedleft
lambda\_opt
\end{minipage} & \begin{minipage}[b]{\linewidth}\raggedleft
CV bal.loss
\end{minipage} & \begin{minipage}[b]{\linewidth}\raggedleft
CV mean SMD
\end{minipage} & \begin{minipage}[b]{\linewidth}\raggedleft
boot mean SMD
\end{minipage} & \begin{minipage}[b]{\linewidth}\raggedleft
CV max SMD
\end{minipage} & \begin{minipage}[b]{\linewidth}\raggedleft
boot max SMD
\end{minipage} \\
\midrule\noalign{}
\endhead
\bottomrule\noalign{}
\endlastfoot
1 & 0.005 & 0.069 & 0.183 & 0.017 & 0.149 & 0.017 \\
2 & 0.011 & 0.068 & 0.183 & 0.017 & 0.144 & 0.017 \\
3 & 0.017 & 0.067 & 0.181 & 0.017 & 0.150 & 0.017 \\
4 & 0.018 & 0.067 & 0.179 & 0.016 & 0.148 & 0.017 \\
5 & 0.032 & 0.067 & 0.181 & 0.016 & 0.145 & 0.017 \\
7 & 0.048 & 0.068 & 0.183 & 0.016 & 0.147 & 0.017 \\
10 & 0.058 & 0.068 & 0.178 & 0.016 & 0.145 & 0.017 \\
12 & 0.067 & 0.068 & 0.184 & 0.016 & 0.147 & 0.017 \\
14 & 0.070 & 0.068 & 0.179 & 0.016 & 0.145 & 0.017 \\
16 & 0.085 & 0.067 & 0.180 & 0.016 & 0.145 & 0.017 \\
18 & 0.092 & 0.068 & 0.180 & 0.016 & 0.146 & 0.017 \\
20 & 0.102 & 0.069 & 0.184 & 0.017 & 0.145 & 0.017 \\
\end{longtable}
}

\end{minipage}%
%
\begin{minipage}[t]{0.50\linewidth}

{\def\LTcaptype{none} % do not increment counter
\begin{longtable}[]{@{}
  >{\raggedleft\arraybackslash}p{(\linewidth - 12\tabcolsep) * \real{0.0875}}
  >{\raggedleft\arraybackslash}p{(\linewidth - 12\tabcolsep) * \real{0.1375}}
  >{\raggedleft\arraybackslash}p{(\linewidth - 12\tabcolsep) * \real{0.1500}}
  >{\raggedleft\arraybackslash}p{(\linewidth - 12\tabcolsep) * \real{0.1500}}
  >{\raggedleft\arraybackslash}p{(\linewidth - 12\tabcolsep) * \real{0.1750}}
  >{\raggedleft\arraybackslash}p{(\linewidth - 12\tabcolsep) * \real{0.1375}}
  >{\raggedleft\arraybackslash}p{(\linewidth - 12\tabcolsep) * \real{0.1625}}@{}}
\toprule\noalign{}
\begin{minipage}[b]{\linewidth}\raggedleft
c\_prop
\end{minipage} & \begin{minipage}[b]{\linewidth}\raggedleft
lambda\_opt
\end{minipage} & \begin{minipage}[b]{\linewidth}\raggedleft
CV bal.loss
\end{minipage} & \begin{minipage}[b]{\linewidth}\raggedleft
CV mean SMD
\end{minipage} & \begin{minipage}[b]{\linewidth}\raggedleft
boot mean SMD
\end{minipage} & \begin{minipage}[b]{\linewidth}\raggedleft
CV max SMD
\end{minipage} & \begin{minipage}[b]{\linewidth}\raggedleft
boot max SMD
\end{minipage} \\
\midrule\noalign{}
\endhead
\bottomrule\noalign{}
\endlastfoot
0.7 & 0.005 & 0.069 & 0.183 & 0.017 & 0.149 & 0.017 \\
1.0 & 0.014 & 0.079 & 0.244 & 0.025 & 0.191 & 0.027 \\
1.5 & 0.026 & 0.105 & 0.326 & 0.239 & 0.241 & 0.058 \\
2.0 & 0.061 & 0.132 & 0.385 & 0.296 & 0.276 & 0.143 \\
2.5 & 0.092 & 0.153 & 0.419 & 0.332 & 0.302 & 0.191 \\
3.0 & 0.116 & 0.169 & 0.445 & 0.352 & 0.315 & 0.211 \\
4.0 & 0.168 & 0.198 & 0.463 & 0.379 & 0.334 & 0.238 \\
5.0 & 0.199 & 0.218 & 0.480 & 0.392 & 0.345 & 0.256 \\
6.0 & 0.185 & 0.229 & 0.488 & 0.398 & 0.346 & 0.266 \\
8.0 & 0.242 & 0.247 & 0.495 & 0.404 & 0.351 & 0.278 \\
16.0 & 0.292 & 0.265 & 0.502 & 0.410 & 0.354 & 0.293 \\
\end{longtable}
}

\end{minipage}%
\newline
\begin{minipage}[t]{0.50\linewidth}

\begin{longtable}[]{@{}
  >{\raggedleft\arraybackslash}p{(\linewidth - 12\tabcolsep) * \real{0.1127}}
  >{\raggedleft\arraybackslash}p{(\linewidth - 12\tabcolsep) * \real{0.1549}}
  >{\raggedleft\arraybackslash}p{(\linewidth - 12\tabcolsep) * \real{0.1690}}
  >{\raggedleft\arraybackslash}p{(\linewidth - 12\tabcolsep) * \real{0.0986}}
  >{\raggedleft\arraybackslash}p{(\linewidth - 12\tabcolsep) * \real{0.1690}}
  >{\raggedleft\arraybackslash}p{(\linewidth - 12\tabcolsep) * \real{0.1268}}
  >{\raggedleft\arraybackslash}p{(\linewidth - 12\tabcolsep) * \real{0.1690}}@{}}
\caption{moderate overlap}\tabularnewline
\toprule\noalign{}
\begin{minipage}[b]{\linewidth}\raggedleft
sigma\_y
\end{minipage} & \begin{minipage}[b]{\linewidth}\raggedleft
lambda\_opt
\end{minipage} & \begin{minipage}[b]{\linewidth}\raggedleft
rel. change
\end{minipage} & \begin{minipage}[b]{\linewidth}\raggedleft
cv.inf
\end{minipage} & \begin{minipage}[b]{\linewidth}\raggedleft
rel. change
\end{minipage} & \begin{minipage}[b]{\linewidth}\raggedleft
boot.inf
\end{minipage} & \begin{minipage}[b]{\linewidth}\raggedleft
rel. change
\end{minipage} \\
\midrule\noalign{}
\endfirsthead
\toprule\noalign{}
\begin{minipage}[b]{\linewidth}\raggedleft
sigma\_y
\end{minipage} & \begin{minipage}[b]{\linewidth}\raggedleft
lambda\_opt
\end{minipage} & \begin{minipage}[b]{\linewidth}\raggedleft
rel. change
\end{minipage} & \begin{minipage}[b]{\linewidth}\raggedleft
cv.inf
\end{minipage} & \begin{minipage}[b]{\linewidth}\raggedleft
rel. change
\end{minipage} & \begin{minipage}[b]{\linewidth}\raggedleft
boot.inf
\end{minipage} & \begin{minipage}[b]{\linewidth}\raggedleft
rel. change
\end{minipage} \\
\midrule\noalign{}
\endhead
\bottomrule\noalign{}
\endlastfoot
1 & 0.026 & & 0.241 & & 0.058 & \\
2 & 0.039 & 1.49 & 0.243 & 1.01 & 0.057 & 0.98 \\
3 & 0.050 & 1.26 & 0.247 & 1.02 & 0.057 & 1.01 \\
4 & 0.057 & 1.15 & 0.242 & 0.98 & 0.056 & 0.99 \\
5 & 0.064 & 1.12 & 0.243 & 1.00 & 0.057 & 1.01 \\
7 & 0.081 & 1.28 & 0.240 & 0.99 & 0.056 & 0.97 \\
10 & 0.134 & 1.65 & 0.243 & 1.01 & 0.056 & 1.01 \\
\end{longtable}

\end{minipage}%
%
\begin{minipage}[t]{0.50\linewidth}

\begin{longtable}[]{@{}
  >{\raggedleft\arraybackslash}p{(\linewidth - 12\tabcolsep) * \real{0.1127}}
  >{\raggedleft\arraybackslash}p{(\linewidth - 12\tabcolsep) * \real{0.1549}}
  >{\raggedleft\arraybackslash}p{(\linewidth - 12\tabcolsep) * \real{0.1690}}
  >{\raggedleft\arraybackslash}p{(\linewidth - 12\tabcolsep) * \real{0.0986}}
  >{\raggedleft\arraybackslash}p{(\linewidth - 12\tabcolsep) * \real{0.1690}}
  >{\raggedleft\arraybackslash}p{(\linewidth - 12\tabcolsep) * \real{0.1268}}
  >{\raggedleft\arraybackslash}p{(\linewidth - 12\tabcolsep) * \real{0.1690}}@{}}
\caption{bad overlap}\tabularnewline
\toprule\noalign{}
\begin{minipage}[b]{\linewidth}\raggedleft
sigma\_y
\end{minipage} & \begin{minipage}[b]{\linewidth}\raggedleft
lambda\_opt
\end{minipage} & \begin{minipage}[b]{\linewidth}\raggedleft
rel. change
\end{minipage} & \begin{minipage}[b]{\linewidth}\raggedleft
cv.inf
\end{minipage} & \begin{minipage}[b]{\linewidth}\raggedleft
rel. change
\end{minipage} & \begin{minipage}[b]{\linewidth}\raggedleft
boot.inf
\end{minipage} & \begin{minipage}[b]{\linewidth}\raggedleft
rel. change
\end{minipage} \\
\midrule\noalign{}
\endfirsthead
\toprule\noalign{}
\begin{minipage}[b]{\linewidth}\raggedleft
sigma\_y
\end{minipage} & \begin{minipage}[b]{\linewidth}\raggedleft
lambda\_opt
\end{minipage} & \begin{minipage}[b]{\linewidth}\raggedleft
rel. change
\end{minipage} & \begin{minipage}[b]{\linewidth}\raggedleft
cv.inf
\end{minipage} & \begin{minipage}[b]{\linewidth}\raggedleft
rel. change
\end{minipage} & \begin{minipage}[b]{\linewidth}\raggedleft
boot.inf
\end{minipage} & \begin{minipage}[b]{\linewidth}\raggedleft
rel. change
\end{minipage} \\
\midrule\noalign{}
\endhead
\bottomrule\noalign{}
\endlastfoot
1 & 0.092 & & 0.302 & & 0.191 & \\
2 & 0.125 & 1.36 & 0.301 & 1.00 & 0.187 & 0.98 \\
3 & 0.090 & 0.72 & 0.302 & 1.00 & 0.190 & 1.01 \\
4 & 0.119 & 1.32 & 0.302 & 1.00 & 0.187 & 0.99 \\
5 & 0.122 & 1.03 & 0.301 & 1.00 & 0.186 & 1.00 \\
7 & 0.173 & 1.42 & 0.301 & 1.00 & 0.188 & 1.01 \\
10 & 0.218 & 1.26 & 0.300 & 1.00 & 0.189 & 1.00 \\
\end{longtable}

\end{minipage}%
\newline
\begin{minipage}[t]{0.50\linewidth}

\begin{longtable}[]{@{}
  >{\raggedleft\arraybackslash}p{(\linewidth - 12\tabcolsep) * \real{0.1000}}
  >{\raggedleft\arraybackslash}p{(\linewidth - 12\tabcolsep) * \real{0.1571}}
  >{\raggedleft\arraybackslash}p{(\linewidth - 12\tabcolsep) * \real{0.1714}}
  >{\raggedleft\arraybackslash}p{(\linewidth - 12\tabcolsep) * \real{0.1000}}
  >{\raggedleft\arraybackslash}p{(\linewidth - 12\tabcolsep) * \real{0.1714}}
  >{\raggedleft\arraybackslash}p{(\linewidth - 12\tabcolsep) * \real{0.1286}}
  >{\raggedleft\arraybackslash}p{(\linewidth - 12\tabcolsep) * \real{0.1714}}@{}}
\caption{overlap axis, sigma\_y = 1}\tabularnewline
\toprule\noalign{}
\begin{minipage}[b]{\linewidth}\raggedleft
c\_prop
\end{minipage} & \begin{minipage}[b]{\linewidth}\raggedleft
lambda\_opt
\end{minipage} & \begin{minipage}[b]{\linewidth}\raggedleft
rel. change
\end{minipage} & \begin{minipage}[b]{\linewidth}\raggedleft
cv.inf
\end{minipage} & \begin{minipage}[b]{\linewidth}\raggedleft
rel. change
\end{minipage} & \begin{minipage}[b]{\linewidth}\raggedleft
boot.inf
\end{minipage} & \begin{minipage}[b]{\linewidth}\raggedleft
rel. change
\end{minipage} \\
\midrule\noalign{}
\endfirsthead
\toprule\noalign{}
\begin{minipage}[b]{\linewidth}\raggedleft
c\_prop
\end{minipage} & \begin{minipage}[b]{\linewidth}\raggedleft
lambda\_opt
\end{minipage} & \begin{minipage}[b]{\linewidth}\raggedleft
rel. change
\end{minipage} & \begin{minipage}[b]{\linewidth}\raggedleft
cv.inf
\end{minipage} & \begin{minipage}[b]{\linewidth}\raggedleft
rel. change
\end{minipage} & \begin{minipage}[b]{\linewidth}\raggedleft
boot.inf
\end{minipage} & \begin{minipage}[b]{\linewidth}\raggedleft
rel. change
\end{minipage} \\
\midrule\noalign{}
\endhead
\bottomrule\noalign{}
\endlastfoot
0.7 & 0.005 & & 0.149 & & 0.017 & \\
1.0 & 0.014 & 2.67 & 0.191 & 1.28 & 0.027 & 1.56 \\
1.5 & 0.026 & 1.96 & 0.241 & 1.26 & 0.058 & 2.15 \\
2.0 & 0.061 & 2.29 & 0.276 & 1.15 & 0.143 & 2.46 \\
2.5 & 0.092 & 1.52 & 0.302 & 1.09 & 0.191 & 1.33 \\
3.0 & 0.116 & 1.26 & 0.315 & 1.04 & 0.211 & 1.10 \\
4.0 & 0.168 & 1.46 & 0.334 & 1.06 & 0.238 & 1.13 \\
5.0 & 0.199 & 1.18 & 0.345 & 1.03 & 0.256 & 1.08 \\
6.0 & 0.185 & 0.93 & 0.346 & 1.00 & 0.266 & 1.04 \\
8.0 & 0.242 & 1.31 & 0.351 & 1.01 & 0.278 & 1.05 \\
16.0 & 0.292 & 1.20 & 0.354 & 1.01 & 0.293 & 1.05 \\
\end{longtable}

\end{minipage}%

\end{figure}%

\hfill\break

\section{}\label{section}

\section{Selector ranking}\label{selector-ranking}

\section{Against the lasso floor}\label{against-the-lasso-floor}

\section{Caveats}\label{caveats}

\section{Appendix}\label{appendix}

\subsection{Appendix 1: Dense outcome, diluted
confounding}\label{appendix-1-dense-outcome-diluted-confounding}

\begin{itemize}
\tightlist
\item
  100 active outcome covariates (\texttt{s\_y\ =\ 100}), 5 confounders.
\item
  \textbf{Confounder signal is diluted:}
  \texttt{beta\_y{[}1:s\_y{]}\ \textless{}-\ 1/sqrt(s\_y)} shrinks every
  outcome coefficient as \texttt{s\_y} grows, including the 5
  confounders.

  \begin{itemize}
  \tightlist
  \item
    Increasing density decreases confounder signal, in effect lowering
    the confounders' share of outcome variance, which \texttt{sigma\_y}
    already does.
  \end{itemize}
\end{itemize}

\includegraphics[width=1\linewidth,height=\textheight,keepaspectratio]{fig_density_sy100.png}

\subsection{Appendix 2: Sparse outcome
replicate}\label{appendix-2-sparse-outcome-replicate}

Same DGP and reps as the sparse figure in Section 2, independent draws.

\includegraphics[width=1\linewidth,height=\textheight,keepaspectratio]{fig_density_fixed_sparse.png}

\subsection{Appendix 3: Sparse vs dense, sigma\_y =
1}\label{appendix-3-sparse-vs-dense-sigma_y-1}

\includegraphics[width=1\linewidth,height=\textheight,keepaspectratio]{fig_density_fixed_sigma1.png}

\subsection{Appendix 4: Noise by overlap, including
awful}\label{appendix-4-noise-by-overlap-including-awful}

\includegraphics[width=1\linewidth,height=\textheight,keepaspectratio]{fig_snr_lasso_all.png}




\end{document}
